\documentclass{article}
\usepackage{graphicx}
\usepackage{svg}

\usepackage[utf8]{inputenc}     % pdfLaTeX
\usepackage[T1]{fontenc}
\usepackage[magyar]{babel}

\title{Kültakaró}
\author{Lengyel Zoltán Péter, Rovenszky Robin}
\date{Legutóbb frissítve: 2026. Március 18.}

\begin{document}

\maketitle
\tableofcontents
\newpage

\section{A bőr}
\begin{figure}
    \centering
    \includegraphics{Assets/Biology/Human/Screenshot 2026-03-25 083130.png}
    \caption{A hámszövet típusai.}
    \label{fig:hámszövet}
\end{figure}

A bőr a testsúlyunk 15-20\%-át teszi ki.
\subsection{Többrétegű laphám}
\begin{itemize}
    \item Alul folyamatosan osztódó sejtek. Ahogy osztódnak tolódnak felfelé, felül elhalt szaruréteget képeznek
\end{itemize}

\subsection{A bőr 3 rétege}
\begin{figure}
    \centering
    \includegraphics{Assets/Biology/Human/Screenshot 2026-03-25 083509.png}
    \caption{A bőr felépítése.}
    \label{fig:bőr_felépítése}
\end{figure}
\begin{itemize}
    \item Hám \begin{itemize}
        \item Nincsenek erek, táplálása az irhából
        \item Szaruréteg
        \item Pigmentsejtek - melanint termelnek, amely véd az UV sugárzástól (melanin, bőrszín) 
    \end{itemize}
    \item Irha \begin{itemize}
        \item Hajszál gyökerei (szőrtüsző) 
        \item Verejtékmirigy
        \item Rengeteg idegvégződés
        \item Erek
    \end{itemize}
    \item Bőralja
    \begin{itemize}
        \item Zsírszövetben gazdag
        \item Összekapcsolja a bőrt az alatta lévő izomcsoportokkal
    \end{itemize}
\end{itemize}

\subsection{Receptorok a bőrben}
\begin{itemize}
    \item Elosztásuk nem egyenletes a testben
    \item Az ingerület agytörzsbe jut
    \item Szabad idegvégződés - Fájdalom - felhívja a figyelmet a veszélyre, és ösztönöz annak elkerülésére (az idegvégződés roncsolódása érzékel)
    \item Tapintó ideg - finom nyomás - ujjak hegyén találhatóak nagy mennyiségben - gyors nyomásváltozásra nagyon érzékenyek
\end{itemize}

\end{document}